\documentclass[10pt,twocolumn,letterpaper]{article}
\usepackage[review]{cvpr}
\usepackage{graphicx}
\usepackage{booktabs}
\usepackage{amsmath,amssymb}
\usepackage{multirow}
\usepackage{xcolor}
\usepackage{url}
\usepackage{microtype}
\graphicspath{{figures/}}

\definecolor{cvprblue}{rgb}{0.21,0.49,0.74}
\usepackage[pagebackref,breaklinks,colorlinks,allcolors=cvprblue]{hyperref}
\def\paperID{*****}
\def\confName{CVPR}
\def\confYear{2026}
\title{Supplementary Material:\
When Agreement Hides Error}

\begin{document}
\maketitle

\section{Scope and Evidence Contract}
This supplement reports protocol details, full strata, and every result that materially limits a submitted claim. Experiments are included according to relevance, not favorability. The manuscript does not claim a new reconstruction architecture, calibrated uncertainty, universal estimator superiority, plane-invariant influence, local failure localization, or prospective acquisition benefit.

The independent units are 24 fastMRI examinations, 115 IXI examinations, and eight NingXia patients containing 22 series. All NingXia intervals resample patients with subordinate series, planes, models, and contexts retained. Dense target slices are used only for retrospective evaluation and never in Anchor-CV, adjacent-anchor, disagreement, or influence scores.

\section{Definitions and Frozen Protocol}
For predictions $f_a=y+e_a$ and $f_b=y+e_b$, shared and observable error fractions are
\begin{equation}
S=\frac{\|\bar e\|^2}{\|\bar e\|^2+\|\delta\|^2},\qquad
O=\frac{\|\delta\|^2}{\|\bar e\|^2+\|\delta\|^2},
\end{equation}
where $\bar e=(e_a+e_b)/2$ and $\delta=(e_a-e_b)/2$. Numerical identity error is below $4.42\times10^{-10}$ in the stored audit.

For nonlinear models, four frozen seeds (6101--6104) generate centered Rademacher fields at 1\% of the globally normalized case range. Each shared anchor is perturbed separately. The response-cosine statistic averages identical-field response cosines across shared anchors and seeds. The broad grid uses reference offsets $(-2,+2)$ and eight target-preserving comparators. Span and asymmetry are explicit covariates; the complete offset list is stored with the experiment registry. The fixed-count subset contains the four comparators with exactly one common anchor.

\section{Controlled Dependence Results}
\begin{table*}[t]
\centering
\caption{Clean same-target dependence control. Stable-wrong $\Delta$ uses high-error/low-disagreement thresholds frozen from the zero-overlap condition.}
\label{tab:clean-full}
\setlength{\tabcolsep}{4pt}
\begin{tabular}{lrrrrr}
\toprule
Cohort & Units & $\rho_{\rm disjoint}$ & $\rho_{\rm shared}$ & $\Delta\rho$ [95\% CI] & Shared / observable error \\
\midrule
fastMRI & 24 examinations & .959 & .867 & .092 [.043,.176] & 97.39\% / 2.61\% \\
IXI & 115 examinations & .981 & .886 & .095 [.078,.126] & 98.42\% / 1.58\% \\
NingXia & 8 patients, 22 series & .892 & .698 & .193 [.111,.371] & 99.75\% / .25\% \\
\bottomrule
\end{tabular}
\end{table*}

The original order-reversal control is not used as the clean causal contrast. It supports the complete covariance chain on IXI and NingXia, but several cosine/shared/observable arrows reverse on fastMRI because order reversal changes more than anchor dependence. The horizontal-flip clean control preserves anchor identities, target, distances, count, $z$ order, mask, and preprocessing and is therefore the primary intervention.

Metric robustness was assessed over MAE, RMSE/MSE, SSIM error, gradient MAE, low-pass MAE, and residual high-frequency MAE for three physical planes, three cohorts, and two architectures. All 42 disjoint-minus-shared point differences are positive. All seven IXI intervals exclude zero. SSIM error has weak or negative absolute ranking on fastMRI/NingXia and several smaller-cohort gradient/high-frequency intervals cross zero. The claim is broad error observability, not metric independence.

\section{Exact Count Falsification and Analytical Influence}
The compact analytical ladder uses all valid sign configurations at fixed absolute distances $\{1,2,3,4\}$. Across fastMRI, IXI, and NingXia, overlap count has overall $+.9$ association with error cosine/shared fraction and $-.9$ to $-1.0$ with observable fraction/disagreement. Every cohort nevertheless shows a zero-to-one reversal. A wider $\{2,4,6,8\}$ ladder reproduces the exception on IXI while NingXia is monotone. This rejects strict count monotonicity without rejecting a broad dependence effect.

\begin{table}[t]
\centering
\caption{Analytical shared influence versus count. Values are Spearman association with error alignment.}
\label{tab:analytic}
\setlength{\tabcolsep}{5pt}
\begin{tabular}{lrrr}
\toprule
Cohort & Count & Common coefficient mass & Gain \\
\midrule
fastMRI & .503 & .815 & .312 \\
IXI & .466 & .789 & .323 \\
NingXia & .630 & .855 & .225 \\
\bottomrule
\end{tabular}
\end{table}

For the four main analytical endpoints, clustered intervals for the absolute-correlation improvement of common coefficient mass over count are all above zero, spanning .189--.283. Signed dot product and coefficient cosine agree qualitatively; common mass is reported because it retains unequal shared magnitude.

\section{Full Neural Shared-Influence Results}
\begin{table*}[t]
\centering
\caption{Broad neural context grid. $\rho_c$ uses overlap count and $\rho_I$ neural shared influence. Advantage intervals are clustered.}
\label{tab:neural-full}
\scriptsize
\setlength{\tabcolsep}{3pt}
\begin{tabular}{llrrrrrr}
\toprule
Cohort & Model & $\rho_c$ align. & $\rho_I$ align. & Adv. [95\% CI] & $\rho_I$ shared & $\rho_I$ observable & partial $\rho_I$ align. \\
\midrule
fastMRI & HAT & .799 & .838 & .039 [.028,.049] & .829 & $-.830$ & .359 \\
 & SwinIR & .796 & .851 & .055 [.048,.061] & .845 & $-.845$ & .395 \\
IXI & HAT & .714 & .797 & .082 [.076,.088] & .794 & $-.795$ & .315 \\
 & SwinIR & .712 & .787 & .075 [.070,.080] & .786 & $-.787$ & .289 \\
NingXia & HAT & .743 & .786 & .043 [.027,.057] & .792 & $-.792$ & .235 \\
 & SwinIR & .735 & .772 & .036 [.016,.057] & .781 & $-.781$ & .197 \\
\bottomrule
\end{tabular}
\end{table*}

The corresponding primary shared-fraction and observable-fraction advantages are .040/.040 and .056/.056 on fastMRI, .089/.089 and .081/.081 on IXI, and .049/.049 and .043/.043 on NingXia for HAT/SwinIR. Every interval excludes zero. Direct disagreement improves in all six strata, but both fastMRI intervals cross zero; it is supporting rather than primary evidence.

\subsection{Fixed-one-overlap test}
\begin{table*}[t]
\centering
\caption{Neural influence partial Spearman at exactly one shared anchor, controlling span, asymmetry, and examination. Error-alignment intervals are shown.}
\label{tab:fixed}
\setlength{\tabcolsep}{4pt}
\begin{tabular}{llrrrr}
\toprule
Cohort & Model & Error alignment & Shared fraction & Observable fraction & Disagreement \\
\midrule
fastMRI & HAT & .169 [.053,.283] & .135 & $-.135$ & $-.161$ \\
 & SwinIR & .350 [.212,.464] & .321 & $-.321$ & $-.286$ \\
IXI & HAT & .346 [.259,.424] & .360 & $-.360$ & $-.543$ \\
 & SwinIR & .258 [.181,.335] & .241 & $-.241$ & $-.587$ \\
NingXia & HAT & .014 [$-.232$,.223] & $-.006$ & .006 & $-.102$ \\
 & SwinIR & $-.039$ [$-.175$,.144] & $-.039$ & .039 & $-.072$ \\
\bottomrule
\end{tabular}
\end{table*}

NingXia does not preserve the expected direction consistently and all intervals cross zero. We therefore do not use wording such as “the direction is preserved but less precise.” The actual result is null/mixed under eight-patient fixed-count inference.

\subsection{Perturbation and plane stability}
The Plan-7 audit varies amplitude (.5/1/2\%), positive/negative field, Rademacher/Gaussian field, and 1/2/4/8 directions. All 42 cohort--architecture--condition strata pass the registered rank-stability threshold; minimum Spearman is .942 and maximum .999. Physical-plane pair stability is .779--.954, and only one of three plane-pair gates passes. This makes the statistic robust to the tested numerical perturbation design but local to a declared physical plane.

\section{Mechanism Chain}
\begin{table*}[t]
\centering
\caption{Partial Spearman ranges across HAT/SwinIR. Local arrows are supported in every stratum. “Ranking” counts architecture intervals excluding zero.}
\label{tab:chain}
\setlength{\tabcolsep}{3.5pt}
\begin{tabular}{lrrrr}
\toprule
Cohort & Influence$\rightarrow$alignment & Alignment$\rightarrow$norm. disagreement & Influence$\rightarrow$observable & Alignment$\rightarrow$ranking \\
\midrule
fastMRI & .372--.410 & $-.933$--$-.926$ & $-.412$--$-.370$ & $-.333$--$-.143$ (0/2) \\
IXI & .321--.347 & $-.893$--$-.890$ & $-.342$--$-.298$ & $-.643$--$-.500$ (2/2) \\
NingXia & .219--.258 & $-.898$--$-.889$ & $-.231$--$-.211$ & $-.667$--$-.548$ (1/2) \\
\bottomrule
\end{tabular}
\end{table*}

The context-ranking endpoint uses only eight configurations and is not a mediation estimate. Influence-to-ranking is still less stable. The manuscript therefore claims that influence is one determinant of local error alignment/observability, not that it alone determines global case ranking.

\section{Anchor-CV, Strong Baselines, and Consequence}
Anchor-CV exactly equals the masked MAE of half the acquired-slice second difference. Frozen-score reproduction has maximum absolute error $5.68\times10^{-14}$. Table~\ref{tab:anchor} reports strong baselines that prevent a uniqueness claim.

\begin{table*}[t]
\centering
\caption{Examination-level $\times4$ difficulty ranking. Lower nAURC is better.}
\label{tab:anchor}
\setlength{\tabcolsep}{5pt}
\begin{tabular}{llrrr}
\toprule
Cohort & Method & Spearman & Kendall & nAURC \\
\midrule
fastMRI & Anchor-CV / adjacent MAE / adjacent RMSE & .943 / .934 / .942 & .812 / .790 / .804 & .0231 / .0222 / .0241 \\
IXI & Anchor-CV / adjacent MAE / adjacent RMSE & .955 / .956 / .956 & .831 / .834 / .826 & .0150 / .0145 / .0147 \\
NingXia & Anchor-CV / adjacent MAE / adjacent RMSE & .958 / .973 / .959 & .844 / .879 / .835 & .0112 / .0072 / .0123 \\
\bottomrule
\end{tabular}
\end{table*}

Partial rank association for Anchor-CV after adjacent MAE, robust scale, and spacing is .333 on fastMRI, $-.059$ on IXI, and $-.100$ on NingXia. Scale normalization reduces/reverses rankings, so the score is not a scale-free invariant. Frozen risk--coverage also has no universal winner: output PC is better on fastMRI/IXI and adjacent MAE is better on NingXia.

For $\times4\!\rightarrow\!\times2$ benefit, Anchor-CV/adjacent MAE/output PC Spearman is .871/.846/.843 on fastMRI and .948/.950/.970 on IXI. All three policies gain above random at every fixed budget. IXI Anchor-CV oracle-normalized gain is .920/.936/.972/.973 at 10/20/30/50\%; output PC is stronger at some budgets. The IXI benefit compares mean HAT/SwinIR $\times4$ error with frozen I3Net $\times2$ error, so it is cross-estimator and retrospective, not a same-model causal treatment effect.

Anchor-CV takes .42/3.97/3.93 ms CPU on representative fastMRI/IXI/NingXia examinations with zero neural calls. Output PC is cheap only after predictions exist. Neither Anchor-CV nor PC localizes the worst physical plane; their supported use is examination triage.

\section{Material Negative and Boundary Results}
\begin{enumerate}
\item Integer overlap is not a monotone dose; the exact analytical ladder reverses at zero-to-one sharing.
\item Fixed-count continuous influence is null/mixed on NingXia; all eight-patient intervals cross zero.
\item Local response rankings are stable to tested perturbation parameters but not invariant across physical planes.
\item The local cancellation chain is stronger than the global context-ranking chain; no causal mediation claim is supported.
\item Anchor-CV is not uniquely informative or universally superior. Adjacent differences are statistically similar and can be stronger; output PC wins several IXI comparisons.
\item Examination scores do not localize the worst plane and are not calibrated pixelwise uncertainty.
\item Densification results are retrospective, cross-estimator on IXI, and omit scanner-time, motion, and clinical utility.
\item Minimum-influence context selection transfers below random/PC on IXI and is heterogeneous on NingXia.
\item Multi-anchor jackknife IQR collapses from development to confirmation/transfer and is far below adjacent MAE.
\item Covariance residualization gives negligible IXI gain with intervals crossing zero and mixed/tied results elsewhere.
\end{enumerate}

Items 8--10 test abandoned mechanism-derived methods. They are not main-paper contributions, but are retained here because they prevent the reader from inferring that the mechanism already yields an improved universal score.

\section{Reproducibility and Evidence Map}
The frozen sources are:
\begin{itemize}
\item clean intervention and stable wrong: \path{CVPR/CVPR_UPGRADE/outputs/clean_max_overlap*} and \path{stable_wrong_prevalence.json};
\item analytical ladder/influence: \path{compact_four_anchor_overlap_ladder*} and \path{shared_influence*};
\item neural influence and chain: \path{plan6_neural_shared_influence_confirm3*} and \path{plan6_mechanism_chain*};
\item Anchor-CV/consequence: \path{plan6_anchor_cv_explanation*} and \path{plan6_ixi_x2_consequence*};
\item Plan-7 stability and negative gates: \path{PLAN7_ACCEPTANCE_OPTIMIZATION/outputs/p7_cvpr_*};
\item full selection rationale: \path{CVPR/CVPR_RECOVERED/PAPER/EVIDENCE_SELECTION_AUDIT.md}.
\end{itemize}

\end{document}
